% Folder 149, batch 03 — archivists' pages 41–60 (author's pp. 20–30).
% Pass: Opus 5.5 (claude-opus-5-5), 2026-10-03 — first pass, unchecked against the pages by a human.
%
% Notation fixed for this batch: E_K for the extension he writes with a
% curly E (the arithmetic fundamental group of K, an extension of Gamma by
% pi_K); Gamma, Gamma_Q for Gal(Qbar/Q); pi_K, pi_U for the geometric
% fundamental groups; B_K, B_U for the étale (Galois) topoi; \mathbf{Q}
% for the rationals. Every page carries mathematics; none skipped.
\documentclass[11pt,a4paper]{article}
\input{../preamble/grothendieck}

\folder{149}
\batch{3}
\pages{41}{60}
\foldertitle{[Autour de La "Longue Marche" à travers la théorie de Galois, pages 1 à 67] : notes manuscrites (s.d.).}
\dating{s.d. — le groupe « Autour de La "Longue Marche" à travers la théorie de Galois » (141 à 149) est daté [à partir de 1978]-1983}
\watermark{Édition de démonstration}

\begin{document}

\page{41}
\note{la page s'ouvre au milieu d'une phrase commencée avant ce lot : il s'agit de reconstruire les plongements de corps à partir des groupes fondamentaux profinis munis de l'opération de $\Gamma$.}
des groupes \add{profinis} extérieurs avec opération de $\Gamma$ dessus. Et on
pressent que le th.~3 du § précédent (« \uncertain{appliqué} » notamment à
$\mathbf{P}^1_K$ convenablement troué …) pourrait donner la clef d'une telle
construction.

Bien \struck{\ill{}} sûr, des homomorphismes extérieurs quelconques
$\pi^{\wedge}_{K'} \to \pi^{\wedge}_{K}$ n'auront pas de sens géométrique — l'idée
est que les opérations du groupe $\Gamma = \Gamma_{\overline{\mathbf{Q}}/k}$
dessus sont si draconiennes, qu'il n'est possible de trouver un hom.\ extérieur
qui y commute que par voie géométrique — par des plongements de
corps.\note{le signe de complétion $\wedge$ est porté au-dessus de l'indice ; sur le second $K$ il est peut-être doublé d'une barre.}
\uncertain{Donc} il est essentiel ici que le corps de base ne soit pas
quelconque, mais un corps tel que $\mathbf{Q}$ (ou, ce qui revient manifestement
au même, une extension de type fini de $\mathbf{Q}$). Encore faut-il se borner
aux hom.\ $\pi_{K'} \to \pi_K$ dont on décrète d'avance que l'image soit
ouverte — sinon,

\page{42}\note{p.~21 de l'auteur.}
prenant pour $\pi_{K'}$ le groupe unité (i.e.\ $K' = k$), on trouverait un hom.\
$K \to k$ correspondant ! Il faut \uncertain{pour} le moins, pour travailler à
l'aise à partir d'homomorphismes $\pi_{K'} \to \pi_K$ (au lieu de
$E_{K'} \to E_K$) supposer que le centralisateur dans $\pi_K$ de l'image de tout
sous-groupe ouvert de $\pi_{K'}$ soit réduit à $(1)$ — on dira que
l'homomorphisme en question est \emph{admissible}. \ill{} de telle façon qu'à
partir des \supplied{hom.} extérieurs (commutant à $\Gamma$) on reconstitue les
classes d'extensions $E_K$ et $E_{K'}$, qui est l'objet vraiment essentiel.
P.~ex.\ si justement $K' = k$, donc $E_{K'} = \Gamma$, ce qui nous intéressera,
ce ne seront pas les $\Gamma$-hom.\ de $\pi_{K'} = (1)$ (!) dans $\pi_K$, mais
bien les \emph{sections} de $E_K$ sur $\Gamma$.

\emph{Question-conjecture.} Soient $K$, $K'$ deux corps, extensions \add{de type}
finies de $\mathbf{Q}$, \struck{considérons} un morphisme $B_{K'} \to B_K$ de
topos sur $B_{\mathbf{Q}}$.

\page{43}
Les conditions suivantes sont-elles bien équivalentes :

a) L'hom.\ provient d'un plongement de corps $K \subset K'$.

b) L'image \add{de l'hom.\ ext.} $E_{K'} \to E_K$ a une image ouverte.

[c) L'hom.\ extérieur $E_{K'} \to E_K$ est admissible.]
\marginal{c) n'est pas \uncertain{juste}, cf.\ plus bas …}

\emph{NB} On sait que a $\Rightarrow$ b $\Rightarrow$ c, et que b) équivaut à
$\pi_{K'} \to \pi_K$ a une image ouverte.\note{l'auteur écrit ici $\pi_K \to \pi_{K'}$ ; le sens des autres flèches de la page est $K' \to K$.}

Une réponse affirmative impliquerait que si
$\operatorname{degtr} K'/\mathbf{Q} < \operatorname{degtr} K/\mathbf{Q}$, il n'y a pas
de tel hom.\ $E_{K'} \to E_K$, compatible avec les projections sur
$E_{\mathbf{Q}} = \Gamma_{\mathbf{Q}}$ ; en particulier, \struck{toute section de}
\add{il en résulterait} que pour toute section de $E_K$ sur
$\Gamma = \operatorname{Im}(E_K \to \Gamma_{\mathbf{Q}})$, ou sur un sous-groupe
\add{$\Gamma'$} ouvert de $\Gamma$, a un centralisateur \add{non}
\uncertain{trivial} dans $E_K$ — et comme son centralisateur dans $\Gamma'$ est
réduit à $(1)$, cela impliquerait (si $\pi_K \neq 1$) que pour une telle
section, on aurait $\pi_K^{\Gamma'} \neq (1)$. Or je n'aperçois que

\page{44}\note{p.~22 de l'auteur.}
ceci est sans doute faux (cf.\ plus bas \add{n°~3}) — il faudrait
\struck{dans} \add{renforcer} c) ci-dessus en

c') L'hom.\ $E^{o}_{K'} \to E_K$ induit par $E_{K'} \to E_K$ est admissible
(où $E^{o}_{K'}$ est le noyau de l'hom.\ composé
\[
  E_{K'} \longrightarrow \Gamma_{\overline{\mathbf{Q}}/\mathbf{Q}}
  \xrightarrow{\ \text{caractère cyclotomique}\ } \widehat{\mathbf{Z}}^{*}\ ).
\]

Mais pour voir que cette condition est \emph{nécessaire} pour que l'hom.\ soit
géométrique, il faudrait \struck{savoir} vérifier que \add{pour} tout sous-groupe
ouvert $E'$ d'un $E_K$, le centralisateur dans $E_K$ (non seulement de $E'$
lui-même, mais même de $E'^{o}$) est réduit à $1$ — ce qui résulte de la
démonstration du th.~1, et du fait* que pour tout sous-groupe ouvert $\Gamma'$
de $\Gamma = \Gamma_{\mathbf{Q}}$, le centralisateur (non seulement de
$\Gamma'$, mais même) de $\Gamma'^{o}$ dans $\Gamma$ est réduit à $(1)$.
\marginal{* à vérifier !}

Donc, la conjecture initiale \uncertain{devrait} être corrigée dans la

\emph{Conséquence (conjecturale).} Pour toute \add{\ill{}} section
\add{d'un sous-groupe} ouvert $\Gamma'$ de $\Gamma_{\mathbf{Q}}$,
\struck{\ill{}} de $E_K$ sur $\Gamma_{\mathbf{Q}}$, \struck{un sous-groupe
ouvert} de sorte que $\Gamma'$ opère (effectivement) sur $\pi_K$, on a (si $K$
pas alg.\ sur $\mathbf{Q}$, i.e.\ $\pi_K \neq (1)$)
$\pi_K^{\Gamma'^{o}} \neq (1)$.

\page{45}
À vrai dire, à certains égards les $\Gamma_K$ sont des groupes trop gros pour
pouvoir travailler directement avec, et il y a lieu de regarder $\Gamma_K$
comme une \emph{limite} de groupes $\Gamma_{U/\mathbf{Q}}$ associés à des
modèles affines de $K$ — et on s'intéressera plus particulièrement à des
modèles affines qui sont des variétés élémentaires — plus généralement, qui
sont des $K(\pi,1)$ (au sens profini …). Il est possible qu'il faille d'ailleurs,
dans l'énoncé de la conjecture, exiger, \struck{\ill{}} pour un hom.\ \add{ext}
$E_{K'} \to E_K$ donné, \ill{} (en plus de l'hyp.\ d'admissibilité et de la
compatibilité avec les hom.\ dans $\Gamma_{\mathbf{Q}}$) qu'il soit compatible
avec les \emph{filtrations} de ces groupes, associées à ces modèles
(« filtrations \uncertain{modéliques} » (grossières)).

Nous allons alors, en même temps que des extensions de t.f.\ de $\mathbf{Q}$,
les hom.\ entre elles, et

\page{46}\note{p.~23 de l'auteur.}
homomorphismes de groupes profinis associés, \struck{regarder} \add{étudier} la
situation analogue pour des « modèles » élémentaires \add{admissibles}, voire
des modèles $K(\pi,1)$ généraux. (On peut aussi regarder de tels modèles sur un
corps \add{$K$ ext.} \struck{\ill{}} de type fini de $\mathbf{Q}$, — mais pour le
moment cette situation mixte me paraît un peu bâtarde …) Si $U$, $V$ sont de
tels modèles, tout morphisme $V \to U$\note{la première lettre est surchargée.}
définit un morphisme de topos galoisiens sur $B_{\mathbf{Q}}$,
$B_U \to B_V$, \struck{\ill{}} et si \ill{} est \struck{admissible} élémentaire
admissible, \ill{} est connu quand on connaît seulement
\[
  H_1(B_{\overline{U}}, \mathbf{Z}_\ell) \longrightarrow
  H_1(B_{\overline{V}}, \mathbf{Z}_\ell)
\]
— ce qui est beaucoup moins que la classe d'iso.\ de hom.\ de
$B_{\mathbf{Q}}$-topos. (En fait, \struck{il suffit} \ill{} \ill{} sur $V$, dès
que $V$ se plonge dans une VA, $f$ est connu quand on connaît sa restriction
sur les topos \ill{} …). Mais quels sont les hom.\ $B_U \to B_V$, ou
$E_U \to E_V$\note{le sens des flèches $B_U \to B_V$, $E_U \to E_V$ est celui de la page, inverse de celui de $V \to U$.}

\page{47}
qui correspondent à des morphismes de modèles ? Avec un peu de culot, on dirait

\emph{Conjecture fondamentale :} Soient $U$, $V$ deux schémas de t.f.\ sur
$\mathbf{Q}$, \struck{\ill{}} \add{\ill{}} réguliers, $U$ une
\struck{modèle} \add{variété} élémentaire admissible sur une extension finie de
$\mathbf{Q}$. \struck{Alors} \add{Considérons} \struck{tout hom.} un hom.\
$B_V \to B_U$ des topos étales sur $\mathbf{Q}$ — ou ce qui revient au même, un
hom.\ de groupes extérieur
$f : E_V = \pi_1(V) \to E_U = \pi_1(U)$, compatible avec les hom.\ \add{ext}
dans $\Gamma_{\mathbf{Q}} = \pi_1(\mathbf{Q})$.
\marginal{NB pour l'unicité, \uncertain{il suffit} \ill{} que $V$ lui-même
soit un modèle élémentaire admissible, \ill{} d'un corps de \ill{} …}

Conditions équivalentes

a) Cet hom.\ provient (à isom.\ près) d'un morphisme $V \to U$ sur les modèles
(qui est alors uniquement déterminé !)

b) $f \mid E_V^{o}$ est admissible, i.e.\ l'image par $f$ de tout sous-groupe
ouvert de $E_V^{o}$ a un centralisateur réduit à $1$

\page{48}\note{p.~24 de l'auteur.}
Pour la nécessité de b), on est ramené aussitôt au cas où $V$ est réduit à un
pt, où cela se réduit à la \struck{conséquence conjecturale de la
conjecture-question préliminaire}.
\marginal{et il \uncertain{vaut} même \uncertain{mieux} se borner à
l'équivalence de a) et b) — la condition c) sur les centralisateurs risque de
\ill{} \ill{} …}\note{aucune condition c) n'est écrite en p.~47 ; la marge vise
peut-être la condition c) des pages 43–44.}

\emph{Conséquence conjecturale.} Soit
$\Gamma' \subset \operatorname{Im}(E_U \to \Gamma_{\mathbf{Q}})$ un sous-groupe
ouvert, correspondant à un corps $k$ fini sur $\mathbf{Q}$, considérons un
$k$-point de $U$, d'où un relèvement $\Gamma' \to E_U$, de sorte que $\Gamma'$
opère sur $\pi_U$. Ceci posé, on a
\[
  \pi_U^{\Gamma'^{o}} = \{1\}.
\]
\note{énoncé encadré d'un trait vertical dans la marge gauche.}

On étudiera par la suite les relations entre cette \add{cons.}
« conjecturale », et la précédente (d'apparence opposée !) concernant les
$E_K$.

La conjecture fondamentale, sur les modèles, implique la conjecture
fondamentale sur les corps, à condition de prendre soin, dans cette dernière,
de se limiter aux hom.\ compatibles aux filtrations modéliques,
\add{\uncertain{importantes}} pour $U$. Plus généralement, prenant des schémas
qui sont des $\varprojlim$ de \struck{schémas} \add{modèles élém.}
\struck{\ill{}} admissibles, \struck{\ill{}} \add{des immersions ouvertes}
morphismes de transition affines

\page{49}
(qui ne proviennent pas : la $\varprojlim$ dans la catégorie des schémas),
\struck{\ill{}} \add{soit} $V$ un schéma $\varprojlim$ de schémas
\uncertain{séparés} réguliers de t.f.\ sur $\mathbf{Q}$ (morphismes de
transition affines \add{immersions ouvertes} \uncertain{sans plus}). Alors les
morphismes \emph{dominants} de schémas $V \to U$ doivent correspondre aux hom.\
ext.\ $E_V \to E_U$, compatibles avec les projections dans
$E_{\mathbf{Q}} = \Gamma_{\mathbf{Q}}$, et tels que l'image soit ouverte.

P.~ex.\ on pourrait prendre (\struck{$U$, $\operatorname{Spec}\mathcal{O}_{X,x}$})
\add{pour} $V$ les spectres d'anneaux locaux \add{réguliers}
$\mathcal{O}_{X,x}$ de \uncertain{schémas} de type fini sur $\mathbf{Q}$.

Cette conjecture fondamentale (éventuellement revue et corrigée en cours de
route !) étant admise, la question qui se pose ensuite est de déterminer les
topos (pro)galoisiens sur $B_{\mathbf{Q}}$ qui proviennent de modèles
élémentaires admissibles — ou encore, les $\pi_U = \pi_1(\overline{U})$ de tels
modèles étant connus, de déterminer quelles sont

\page{50}\note{p.~25 de l'auteur.}
\supplied{les opér}ations extérieures possibles de sous-groupes \add{ouverts}
$\Gamma'$ de $\Gamma_{\mathbf{Q}}$ sur de tels groupes fondamentaux — et
éventuellement question analogue pour d'autres types de groupes profinis,
correspondant à des $K(\pi,1)$ qui se réaliseraient par des variétés
algébriques (sur $\mathbf{C}$, disons), mais pas par des variétés élémentaires.
(J'ai en vue \struck{surtout} des \struck{topos} \add{variétés} modulaires,
telles que, notamment, des variétés modulaires pour les courbes algébriques …)
À partir de là, on reconstruirait par recollement, en termes profinis, tous les
schémas lisses sur un corps de type fini sur $\mathbf{Q}$ (\uncertain{ou
plutôt}, la catégorie de \uncertain{ceux-ci} …), et plus gén.\ sur un corps
quelconque — puis, sans doute, par « recollement », la catégorie des schémas
loc.\ de type fini sur un $K$ — \uncertain{ou du} moins la cat.\ des
\struck{\ill{}} \add{fractions} qui s'en déduit en rendant inv.\ les

\page{51}
homéomorphismes universels …

Les réflexions précédentes suggèrent aussi des énoncés comme le suivant : Pour
un schéma de base $S$ loc.\ noeth.\ donné, le foncteur $X \mapsto X_{\text{et}}$,
allant de la catégorie des schémas \add{réduits} \struck{réduits} loc.\ de prés.\
finie sur $S$, vers la 2-catégorie des topos au-dessus de $S_{\text{et}}$, est
\struck{pleinement} 1-fidèle (donc hom.\ $f, g : X \rightrightarrows Y$ tels que
les morphismes de topos $X_{\text{et}} \rightrightarrows Y_{\text{et}}$
\add{$f_{\text{et}}, g_{\text{et}}$} au-dessus de $S_{\text{et}}$ soient
isomorphes, sont égaux) et est peut-être \emph{pleinement fidèle}, quand on
passe à la catégorie de fractions de $(\mathrm{Sch}_{\text{l.t.f}})/S$ obtenue
en rendant inv.\ les homéom.\ universels …
\marginal{$S$ de car.~$0$ ?}
Exprimant \struck{\ill{}} ceci p.~ex.\ par les automorphismes d'une courbe alg.\
propre sur une ext.\ finie de $\mathbf{Q}$, on retrouverait le « fait » que tout
automorphisme ext.\ de $E_K$ ($K$ le corps des fonctions de $X$) qui respecte la
structure : \ill{} et qui commute à l'action de $\Gamma$, provient d'un
automorphisme de $X$.

\page{52}\note{p.~26 de l'auteur.}

\section*{3. Étude des sections de $E_U$ sur $\Gamma$.}

Soit $U$ un \struck{variété élémentaire sur} \add{schéma \uncertain{géom.\
connexe} lisse de t.f.} le corps $K$, d'où $E_U \to E_K$, et on se propose
d'étudier les sections mod $\pi_{U/K}$-conjugaison — plus généralement, en même
temps, regarder les sections $E'_K \to E_U$, où $E'_K$ est un sous-groupe ouvert
de $E_K$ (ce qui signifie qu'on fait une extension de base $K'$ de type fini sur
le corps $\mathbf{Q}$ et finie sur $K$).
\marginal{On a choisi un rev.\ universel $\widetilde{U}$ de $U$ pour définir
$\overline{K}$, et $E_U$, $E_K$, et $E_U \to E_K$.}
\[
  (1)\qquad 1 \longrightarrow \pi_{U/K} \longrightarrow E_U \longrightarrow
  E_K \longrightarrow 1 .
\]
\note{la suite (1) est écrite verticalement dans la marge gauche, à côté de ce passage ; les renvois à (1) des pages suivantes la visent.}
Si $U$ se plonge dans un schéma en groupes commutatif \struck{\ill{}} \add{\ill{}}
rigides, \struck{\ill{}} l'application
\[
  U(K) \longrightarrow
  \begin{array}{l}
    \text{classes d'isomorphie de sections de } B_U \text{ sur } B_K,\\
    \text{i.e.\ classes de } \pi_{\overline{U}/K}\text{-conjugaison}\\
    \text{de sections de } E_U \text{ sur } E_K
  \end{array}
\]
est injective. On va examiner d'autres façons \add{« géométriques »} de trouver
des sections.

Supposons d'abord que $U$ est une courbe algébrique, \struck{qu'il ne soit pas
\ill{}} \add{\ill{}} type $(0,0)$ ou $(0,1)$ i.e.\
$\pi_1(\overline{U}) = \pi_{U/K} \neq \{0\}$ ; on a vu que pour \ill{}
$i \in \widehat{U} \setminus U$ (« pts à l'infini », \ill{}), le groupe de
lacets $L_i$ fournit un \ill{} (\ill{} \ill{} si $i$ rat./$K$, i.e.\ fixé

\page{53}
par $E_{\overline{K},K}$) en prenant son centralisateur $Z(L_i)$ dans $E$, d'où
\[
  (2)\qquad 1 \longrightarrow L_i \longrightarrow Z(L_i) \longrightarrow
  \Gamma \longrightarrow 1,
  \qquad L_i \simeq T_\infty(\overline{K}^{*}) = T,
\]
et on prend les scindages de cette extension. Il en existe, p.~ex.\ définis par
une uniformisante de $\mathcal{O}_{\widehat{U},i}$. L'ens.\ des classes de
conjugaison \struck{\ill{}} des scindages de (2) est un torseur sous
\struck{$H^1(\Gamma,T)$}
\[
  (3)\qquad H^1(\Gamma, T) \simeq \varprojlim_n H^1(K, {}_n\mu) \simeq
  \widehat{K^{*}} = \varprojlim_n K^{*}/K^{*n}
\]
et \struck{\ill{}} s'envoie injectivement dans l'ens.\ des classes de
$\pi$-conj.\ de scindages.\note{le numéro (3) servira de nouveau en p.~54 pour une autre relation.}

\emph{Proposition.} On suppose $(g,\nu) \neq (0,0), (0,1)$ i.e.\
$\pi_{-} = \pi_{U/K} \neq (1)$. Alors les classes de $\pi$-conjugaison
\struck{associées} \add{de} scindages de (1) définis par les scindages de (2)
sont distinctes de celles associées aux pts de $U(K)$. \struck{De plus},
\add{(de plus,} si \ill{} dans le cas anabélien, $(g,\nu) \neq (0,2)$,
\struck{\ill{}} \add{i.e.\ si \ill{}}) alors les classes de
$\pi$-conjugaison de scindages de (1), associés à des scindages de (2) pour deux
indices $i = i_1$ et $i = i_2$ distincts, sont distinctes.
\marginal{C'est démontré sauf peut-être pour le type $(0,3)$, qui mérite un
examen approfondi.}

\page{54}\note{p.~27 de l'auteur.}
La première assertion s'obtient en « bouchant » le trou $i$, donc la section
envisagée devient la section de $U \cup \{i\} = U'$ associée au pt $i$, et elle
est donc distincte de celles associées aux \struck{pt} autres pts de $U'$, i.e.\
aux pts de $U$ — à fortiori aussi \ill{} \uncertain{quotient} par le s-groupe
inv.\ \uncertain{engendré} par $L_i$. On raisonne de même pour voir que les
classes de conj.\ de scindages associés à un $L_{i_1}$ et un $L_{i_2}$,
$i_1 \neq i_2$, sont distincts — \struck{\ill{}} \uncertain{au moins} dans tous
les cas sauf le cas du type $(0,3)$, — en bouchant deux trous, on tombe sur le
type $(0,1)$, où on ne dispose pas du résultat d'injectivité. Mais on peut s'en
tirer autrement, à condition d'établir que pour un scindage de (2), faisant
opérer $\Gamma$ sur $\pi$, on a
\[
  (3)\qquad \pi^{\Gamma^{o}} = L_i
\]
(donc $\pi^{\Gamma} = (1)$, d'ailleurs) — résultat qui me paraît hautement
plausible.
\marginal{À énoncer parmi les conjectures-clefs …}
Cela

\page{55}
montrerait que la classe de conjugaison de section détermine la classe de
conjugaison de $L_i$, donc $i$.

\emph{Conjecture A.} Soit $U$ courbe alg.\ anabélienne géom.\ connexe sur un
corps $K$ de type fini sur $\mathbf{Q}$. Alors toute section de (1) est
\struck{définie} d'un des deux types précédents, i.e.\ soit définie par un pt de
$U(K)$, soit par une section d'une extension (2), avec
$i \in I(\overline{K})^{\Gamma}$, i.e.\ $i$ un pt de
$\widehat{U} \setminus U$, rationnel sur $K$.
\marginal{Il y a plus : des précisions : $\pi^{\Gamma^{o}} = (1)$ (premier cas),
$\pi^{\Gamma^{o}} = L_i$ (deux.\ cas).}

Si on admet cette conjecture, alors on en conclurait, par passage à la limite,
en considérant le corps des fonctions $L$ de $U$ et $E_L \to E_K$ ($E_L$ pro-\ill{}
considéré comme un groupe à lacets « infini » (avec une infinité de \add{classes
de} sous-groupes $L_i$ de lacets …)) que tout scindage de cette extension provient
d'un scindage d'une extension du type (2), avec
$i \in I^{\Gamma} = X(K)$ $(X = \widehat{U})$. Les classes \add{de} tels
scindages

\page{56}\note{p.~28 de l'auteur.}
se groupent donc \struck{\ill{}} \add{classes de} paquets (en regardant les
centralisateurs des s-groupes images de $\Gamma^{o}$ par ces sections), et on
\ill{} mettre ensemble les scindages qui ont des
$\operatorname{Centr}(\Gamma^{o})$ conjugués (même s'ils ne sont pas eux-mêmes
conjugués). Donc on retrouverait une description de $X(K)$ (comme aussi des
$X(K')$ pour $K'$ ext.\ finie de $K$) en termes de l'extension $E_L$ de $E_K$ par
$\pi_{L/K}$, en même temps qu'une façon de reconstituer les $U = X \setminus I$ …

Donc en fait c'est la structure $E_L \to E_K$ qui est la plus riche a priori, et
de loin plus commode pour le genre $0$ et $1$, vu que la considération des $U$ de
type $(g,\nu)$ $(2g+\nu \geq 3)$ casse le groupe « continu »
d'automorphismes … La forme « modélique » de la conjecture précédente
\uncertain{revient} à la forme « birationnelle », quand on y précise cette
dernière en disant que

\page{57}
les scindages de $E_U \to E_K$ se \struck{\ill{}} \uncertain{remontent} en un
scindage de $E_L \to E_K$ (ou aussi, pour un scindage de $E_V \to E_K$, si $V$
est un modèle $\subset U$).

On va mettre les conjectures précédentes (sous forme modélique, disons) sous une
forme plus géométrique, en introduisant, en même temps que un revêtement
universel $\widetilde{U}$ de $U$, la normalisation $X'$ \add{\ill{}} de $X$
\add{dans $\widetilde{U}$} (où $X = \widehat{U}$). (NB je m'abstiens de la noter
$\widetilde{X}$, vu que justement il n'est pas étale sur $X$ …). Notons que pour
\add{l'un des} $i \in I = \overline{X} - \overline{U}$, les $L_i$ dans
$\overline{\pi} = \pi(\overline{U})$ sont en corr.\ 1-1 avec \struck{l'ens.} les
fibres $X'_i$ de $X'$ au-dessus de $i$

\begin{tikzcd}
X \arrow[r, no head] & \overline{X} \arrow[r, no head] & X' \\
U \arrow[r, no head] \arrow[u, no head, "\cup" description] \arrow[d, no head]
 & \overline{U} \arrow[r, no head] \arrow[u, no head, "\cup" description] \arrow[d, no head]
 & \widetilde{U} \arrow[u, no head, "\cup" description] \\
K \arrow[r, no head] & \overline{K} &
\end{tikzcd}

Donc $X'$ peut être considéré comme la réunion de $\widetilde{U}$, et de
$X' \mid I$ = ensemble des sous-groupes lacets de $\overline{\pi}$, qui
\uncertain{joueraient} ainsi comme des « pts à l'infini » du rev.\ universel
$\widetilde{U}$. D'ailleurs $E_U$ s'interprète

\page{58}\note{p.~29 de l'auteur.}
comme le groupe des \struck{$K$}-automorphismes du schéma $\widetilde{U}$
\add{compatibles avec}, et $E_U \to E_K$ comme l'hom.\ de passage au quotient
évident (\emph{NB} $\overline{K}$ s'identifie à la clôture alg.\ de $K$ dans
$H^0(\widetilde{U}, \mathcal{O}_{\widetilde{U}})$, donc $E_U$ opère sur
$\operatorname{Spec}\overline{K}$ de façon naturelle …). Une section de $E_U$
sur $E_K$ est donc une action de $E_K$ sur $\widetilde{U}$, compatible avec son
action sur \uncertain{$\overline{K}$}, et continue dans un sens convenable
(sans doute la conv.\ simple-discrète sur les \struck{\ill{}} $U_i$ finis sur
$\overline{U}$ entre $\overline{U}$ et $\widetilde{U}$ …).

Considérons donc la

\emph{Conjecture B.} Toute telle action de $\Gamma$ sur $\widetilde{U}$ admet
\emph{dans $X' = \widehat{\widetilde{U}}$} un pt fixe et un seul.
\note{énoncé encadré d'un trait vertical dans la marge gauche.}
\marginal{et l'action induite de $\Gamma^{o}$ doit avoir un pt fixe et un seul
— cf.\ plus bas.}

Ceci signifie donc

a) S'il y a un pt fixe à distance finie, i.e.\
$\tilde{x} \in \widetilde{U}^{\Gamma}$, alors 1°) L'image de $\tilde{x}$ dans
$U$ est uniquement déterminée — c'est essentiellement le th.~3 du §~1 (deux pts
distincts de $U(K)$ définissent des classes de conj.\ de scindages distinctes) et

\page{59}
2°) $\pi^{\Gamma} = (1)$ (i.e.\ il n'y a pas d'autre pt fixe dans
$\widetilde{U}$ au-dessus de $x \in U(K)$), et enfin

3°) il n'y a pas en même temps un pt fixe \struck{\ill{}} $\tilde{x}$ à
l'infini — i.e.\ il n'existe pas de $L_i$ normalisé par $\Gamma$, i.e.\ un
scindage du premier type n'est pas en même temps du deuxième (fait que nous
avons établi directement, précédemment).
\marginal{C'est un cas particulier d'une \uncertain{conjecture} « conjecturale »
des conjectures \ill{} ; il faudrait \uncertain{écrire}
$\pi^{\Gamma^{o}} = (1)$.}

D'autre part, dans le cas de pts fixes à l'$\infty$, l'unicité de l'image dans
$X$ signifie qu'une même action effective ne peut normaliser à la fois un $L_i$
et un $L_j$ $(i \neq j)$ — fait également établi sauf dans le cas
$(g,\nu) = (0,3)$ — et l'unicité au-dessus d'un $i \in I$ fixé \add{($i$ fixé)}
signifie que le $L_i$ normalisé par $\Gamma$ est unique, ce qui est un
affaiblissement de la relation
\[
  L_i = \pi^{\Gamma^{o}}
\]
pour ces opérations, conjecturée plus haut.

\page{60}\note{p.~30 de l'auteur.}
En fait, je conjecture que dans la conj.~B, il est même vrai que $\Gamma^{o}$
agissant sur $X' = \widehat{\widetilde{U}}$ a un pt fixe et un seul (ce qui est
plus fort, \uncertain{car} ce \add{fixe} serait nécessairement fixe par
$\Gamma$). Ceci \struck{signifie} \add{implique}, dans le cas de pts fixes à
distance finie, que $\pi^{\Gamma^{o}} = (1)$, comme il se devrait en général —
et dans le cas de pt fixe à l'infini, que
\[
  \pi^{\Gamma^{o}} \subset \operatorname{Norm}_{\pi}(L_i) = L_i
\]
donc la relation $\pi^{\Gamma^{o}} = L_i$ « hautement plausible » de tantôt !

Du coup je me demande d'ailleurs si ces énoncés ne devraient pas être
\uncertain{valables} sans partir d'un relèvement de $\Gamma$ tout entier, mais
seulement de $\Gamma^{o}$ (ou d'un sous-groupe contenant un sous-groupe ouvert
de $\Gamma^{o}$, plus gén.) — ce qui signifierait plus ou moins que
\note{la phrase continue au-delà de ce lot.}

\end{document}
